\subsection{Simultaneous pole contributions to the Fej\'er energy}
\label{ssec:fejer-multipole}

The signed Fej\'er energy also gives quantitative information at a
fixed exponent. The following bound records the contributions of any
finite set of poles on the same vertical line.

\begin{proposition}[A sum of pole contributions]\label{prop:fejer-multipole}
Put
\[
 M(s)=\frac{\zeta(2s)}{s\zeta(s)},\qquad
 K(s)=1+\frac{2}{1+2\R s}-2\R\frac1{1+s}.
\]
Fix \(\beta\ge1/2\) and an integer \(r\ge1\). For a finite set
\(\mathcal P\) of distinct poles \(\rho=\beta+i\gamma\) of \(M\)
of order exactly \(r\), write
\(A_\rho=\lim_{s\to\rho}(s-\rho)^rM(s)\). Then
\begin{equation}\label{eq:fejer-multipole}
 \limsup_{N\to\infty}
  \frac{Q_\lambda(N)}{N^{2\beta}(\log N)^{2r-2}}
 \ge \sum_{\rho\in\mathcal P}K(\rho)
              \left|\frac{A_\rho}{(r-1)!}\right|^2.
\end{equation}
If the limit superior is finite, \(M\) has no pole in
\(\R s>\beta\) and every pole on \(\R s=\beta\) has order at
most \(r\).
\end{proposition}

\begin{proof}
Take \(L(X)=0\) on \(0<X<1\). In the complex Hilbert space
\(\mathcal H=L^2((0,1);\mathbb C^2)\), put
\[
 U_X(u)=(L(uX),L(X)-L(uX)),\qquad w_s(u)=(u^s,1-u^s).
\]
For an integer \(N\), the value \(L(uN)=L(j)\) occupies the
interval \([j/N,(j+1)/N)\) for \(0\le j<N\). By
\eqref{eq:fejerboundary}, this gives
\(\|U_N\|_{\mathcal H}^2=Q_\lambda(N)\). For \(X=N+a\),
\(0\le a<1\), the substitution \(v=uX\) gives the exact
interpolation
\begin{equation}\label{eq:fejer-norm-interpolation}
 \|U_X\|_{\mathcal H}^2
   =\frac{NQ_\lambda(N)+a|L(N)|^2}{X}.
\end{equation}
Since \(|L(N)|^2\le2Q_\lambda(N)\) by \eqref{eq:completion},
the numerator in \eqref{eq:fejer-multipole} can be replaced by
\(\|U_X\|_{\mathcal H}^2\) and the limit superior taken over
real \(X\), with the denominator at \(X\). Denote the resulting
limit superior by \(C\).

For \(\R s>1\), \(\|U_X\|_{\mathcal H}\ll X\) ensures absolute
convergence of the Hilbert space integral
\begin{equation}\label{eq:fejer-vector-mellin}
 \int_1^\infty U_X X^{-s}\frac{\dd X}{X}=M(s)w_s.
\end{equation}
For the first coordinate, substitute \(v=uX\), and use
\(L(v)=0\) below~1. For the second, subtract the first coordinate
from the transform of \(L(X)\). The scalar transform of \(L\) is
\(M(s)\) by partial summation and the Euler product. The map
\(s\mapsto w_s\) is holomorphic into \(\mathcal H\) on
\(\R s>-1/2\), since on compact subsets the derivatives of
\(u^s\) are dominated in \(L^2\) by a power of \(u\) times a
power of \(|\log u|\). Direct integration gives
\(\|w_s\|_{\mathcal H}^2=K(s)>0\) there.

If \(C=\infty\), \eqref{eq:fejer-multipole} is immediate.
Otherwise put \(F(t)=e^{-\beta t}U_{e^t}\) for \(t\ge0\).
The definition of \(C\) gives
\(\|F(t)\|_{\mathcal H}\ll 1+t^{r-1}\).
Consequently the integral in \eqref{eq:fejer-vector-mellin}
is holomorphic for \(\R s>\beta\). Meromorphic continuation
identifies it with the displayed right side there. As
\(w_s\ne0\), this excludes poles of \(M\) in that open
half-plane. At a boundary point \(\rho\), the same bound on
\(F\) gives \(\|M(\rho+\delta)w_{\rho+\delta}\|_{\mathcal H}
=O(\delta^{-r})\) as \(\delta\downarrow0\), which excludes
a pole of order greater than \(r\).

For each \(\delta>0\), equip the space of \(\mathcal H\)-valued
functions on \([0,\infty)\) with the \(L^2\) norm for the
measure \(e^{-\delta t}\dd t\). Denote its inner product by
\(\langle\cdot,\cdot\rangle_\delta\), linear in its first
argument, and put \(G_r=\Gamma(2r-1)\). For every
\(\varepsilon>0\), the bound
\(\|F(t)\|_{\mathcal H}^2\le(C+\varepsilon)t^{2r-2}\) holds
for all sufficiently large \(t\). Its integral over a fixed
bounded interval remains bounded as \(\delta\downarrow0\).
Thus Euler's gamma integral gives
\begin{equation}\label{eq:fejer-multipole-abel}
 \limsup_{\delta\downarrow0}\delta^{2r-1}\|F\|_\delta^2
 \le C G_r.
\end{equation}

For each \(\rho=\beta+i\gamma\in\mathcal P\), take the test
function
\[
 g_{\rho,\delta}(t)=\frac{\delta^{r-1/2}}{\sqrt{G_r}}
    t^{r-1} e^{i\gamma t}\,v_\rho,
 \qquad v_\rho=\frac{w_\rho}{\|w_\rho\|_{\mathcal H}}.
\]
Its norm is one. For \(\rho\ne\rho'\), the magnitude of
\(\langle g_{\rho,\delta},g_{\rho',\delta}\rangle_\delta\)
is at most
\((\delta/|\delta-i(\gamma-\gamma')|)^{2r-1}\), again by
the gamma integral. Hence the largest eigenvalue of this finite
Gram matrix tends to one. By \eqref{eq:fejer-vector-mellin},
\[
 \langle F,g_{\rho,\delta}\rangle_\delta
  =\frac{\delta^{r-1/2}}{\sqrt{G_r}}
  \left\langle
   (-1)^{r-1}\frac{\dd^{r-1}}{\dd s^{r-1}}
           [M(s)w_s]\bigg|_{s=\rho+\delta},v_\rho
  \right\rangle_{\mathcal H}.
\]
The signed derivative of the leading Laurent term
\(A_\rho w_\rho(s-\rho)^{-r}\) is
\[G_r A_\rho w_\rho(s-\rho)^{-(2r-1)}/(r-1)!.\]
Consequently
\[
 \lim_{\delta\downarrow0}\delta^{r-1/2}
       \langle F,g_{\rho,\delta}\rangle_\delta
   =\sqrt{G_r}\,\frac{A_\rho}{(r-1)!}\|w_\rho\|_{\mathcal H}.
\]
The sum of the squared scalar products with the tests is at
most \(\|F\|_\delta^2\) times their largest Gram eigenvalue.
Indeed the operator sending coordinate vectors to the tests
has squared norm equal to that eigenvalue, and the analysis
map taking scalar products is its adjoint. Multiply by
\(\delta^{2r-1}\), let \(\delta\downarrow0\) using
\eqref{eq:fejer-multipole-abel}, and divide by \(G_r\).
This proves \eqref{eq:fejer-multipole}.
\end{proof}

\begin{corollary}[A critical-line lower bound]\label{cor:fejer-critical-poles}
For every finite set \(\mathcal Z\) of simple critical-line
zeros of \(\zeta\),
\begin{equation}\label{eq:fejer-critical-poles}
 \limsup_{N\to\infty}\frac{Q_\lambda(N)}N
 \ge\frac{2}{3|\zeta(1/2)|^2}
 +\sum_{\rho\in\mathcal Z}
   \left(2-2\R\frac1{1+\rho}\right)
   \left|\frac{\zeta(2\rho)}{\rho\zeta'(\rho)}\right|^2.
\end{equation}
The left side is infinite if RH fails or if any critical-line
zero has multiplicity greater than one.
\end{corollary}

\begin{proof}
At \(s=1/2\) the residue of \(M\) is \(1/\zeta(1/2)\),
and \(K(1/2)=2/3\). At a simple critical zero \(\rho\), its
residue is \(\zeta(2\rho)/(\rho\zeta'(\rho))\). The numerator
is nonzero by the classical nonvanishing of \(\zeta\) on
\(\R s=1\). Apply Proposition~\ref{prop:fejer-multipole} with
\(\beta=1/2\), \(r=1\) and \(\mathcal P=\{1/2\}\cup\mathcal Z\).
Any right-hand zero, or critical zero of higher multiplicity,
would give an excluded pole in the case of a finite limit
superior. The functional equation reflects any off-line zero
to the right of the critical line if necessary.
\end{proof}

The verified simple zeros of \cite{pt} give strictly positive
terms in \eqref{eq:fejer-critical-poles}. Each finite set of
such zeros gives a lower bound, and adding further zeros
increases it. The estimate concerns the lower side of the
energy question. The square-root-scale upper bound with an
arbitrary power loss remains open.
